% Prospective R124B inference, fixed after A and before B source/state freezes.
\section{The prospective second experiment and joint A--B coverage}
\label{supp:inference124b}

The second design is selected after the first experiment has been
observed. Its conclusions therefore concern a newly frozen service and
new independent states, conditional on that preparation. They do not
relabel the first experiment's design as having anticipated its outcome.
Let $\mathcal H_A$ contain all first-experiment information and all
subsequent preparation available before the second draw. Conditional on
$\mathcal H_A$, the scientific source, actual actors, service definitions,
candidate menu, reporting families and all numerical tolerances are fixed.
The 32 second-experiment initial states are independent draws from the
declared two-coordinate 40-bit dyadic law, independent of
$\mathcal H_A$. This premise concerns the initial-state population;
innovations inside the original Bellman operators retain their continuous
law. Timing repeats, multiple risk specifications and different actors on
the same initial state are not additional independent units.

All statements below are conditional on sound numerical enclosures and
the stated entropy assumption. Each realized service must be the fixed
measurable rule declared before the second draw. In particular, a failed
certificate is handled by its frozen retention rule, not by deleting the
state or replacing the actor after observing its result.

\begin{proposition}[The second experiment's complete inference family]
\label{prop:inference124b}
Put $n_B=32$, $\tau=1/256$ and $t=2^{-24}$. There are 58 fixed
actor/risk cells and three services per cell: the completed fused service,
the paid-value bridge service, and the joint service. Define 232
nonnegative actual original-policy gains as follows: each of the 174
services has a gain relative to its own qualified stage incumbent, and
each of the 58 joint services also has its cumulative gain relative to
the completed fused service. The stage incumbents are, respectively,
the original qualified incumbent before fused enhancement, the completed
fused policy, and the completed bridge policy. All these gains belong to
$[0,\tau]$.

For each gain, let $[\bar\ell,\bar u]$ be a sound numerical enclosure
of its sample mean. The interval
\begin{equation}\label{eq:nonnegativeinference124b}
 \left[\max\{0,\bar\ell-7\tau/16\},
       \min\{\tau,\bar u+7\tau/16\}\right]
\end{equation}
contains its population mean simultaneously across all 232 gains except
on an event of conditional probability at most
$464\exp(-49/4)<1/100$.

For each of the same 232 gains, count the states on which the fixed
exact certificate guarantees a gain at least $t$. At count $k$, use
the lower one-sided binomial probability bound with per-event failure
$\alpha_B=1/23200$. Let its conservative numerical endpoint be
$p^-_k$, including $p^-_0=0$. Then $t p^-_k$ is a population mean-gain
lower bound simultaneously across this family except on an event of
conditional probability at most $1/100$.

The 348 signed comparisons are all 174 services against the same-risk
online/fused service and all 174 against the same-budget
quadratic/VALUE/fused service, with budget 5 fixed for the online cell.
Positive means that the service being evaluated has the lower original
cost. Each signed difference belongs to $[-\tau,\tau]$. Its sound
sample-mean interval enlarged in both directions by $7\tau/8$ and
intersected with $[-\tau,\tau]$ covers the population difference
simultaneously across all 348 comparisons except on an event of
conditional probability at most $696\exp(-49/4)<3/100$.

Consequently, all these intervals and lower bounds hold jointly with
conditional coverage at least $19/20$. A tighter bound on the same
failure event is
\begin{equation}\label{eq:actualfailure124b}
 \varepsilon_B=1160\exp(-49/4)+1/100
              <79/5000=0.0158.
\end{equation}
The interval for a mathematically identical service compared with itself
is exactly zero. Equality observed only on the sampled states establishes
a sample identity, not a population identity.
\end{proposition}

\begin{proof}
The starting incumbent and all completed fused, bridge and joint policies
are qualified under the same original objective, with cost between
$V_0$ and $V_0+\tau$ at every state. The fused acceptance certificate
and Theorem~\ref{thm:stagedpolicy124b} make each stage gain nonnegative.
Their upper bound is $\tau$, because the origin costs no more than
$V_0+\tau$ and the new policy cannot cost less than $V_0$.
The two enhancement gains telescope, so the joint-versus-fused gain
also belongs to $[0,\tau]$, rather than merely to $[0,2\tau]$.
The sum of the two proved stage gain lower bounds is a lower bound
on this cumulative gain. Retention has exactly zero gain. Two qualified
policies for the same risk and state are in the same interval of length
$\tau$; subtracting proves the support for every signed comparison.

Conditional on $\mathcal H_A$, the observations within each fixed cell
are independent and have the same law. Hoeffding's two-sided inequality
for an interval of length $R$ gives failure at most
$2\exp(-2n_B r^2/R^2)$ at radius $r$. For
$R=\tau$, $r=7\tau/16$, and for $R=2\tau$,
$r=7\tau/8$, the exponent in both cases is $49/4$.
The two family union bounds therefore have factors $2(232)=464$
and $2(348)=696$. Exact rational arithmetic verifies
\begin{equation}\label{eq:exponentialcheck124b}
 \sum_{j=0}^{30}\frac{(49/4)^j}{j!}>200000.
\end{equation}
All omitted exponential-series terms are positive. Consequently
\begin{align*}
 464\exp(-49/4)&<464/200000<1/100,\\
 696\exp(-49/4)&<696/200000<3/100.
\end{align*}
Dependence across methods
on a shared state does not enter these union bounds. Replacing the exact
sample mean by a sound enclosing interval and clipping to the proved
support preserves coverage.

For the certificate event, let $p$ denote its conditional population
probability. Conditional on $\mathcal H_A$, its count has law
$\operatorname{Bin}(32,p)$, since the certificate and threshold are fixed
before the draw. For $k>0$ the exact lower endpoint solves
\[
 T_k(p):=\sum_{j=k}^{32}\binom{32}{j}p^j(1-p)^{32-j}
             =\alpha_B.
\]
This polynomial is continuous and strictly increasing on $(0,1)$,
with endpoint values zero and one. Binomial tail inversion gives
$\Pr\{p<p^-_K\mid\mathcal H_A\}\le\alpha_B$ when the exact
endpoint is used. Equivalently, the violating counts form an upper
tail of a binomial law whose probability is at most $\alpha_B$;
the inequality at equality and discrete attainable counts can only
make the bound conservative. Starting with the bracket $[0,1]$, the
specified 48 rational bisections retain the lower endpoint of a bracket
containing the root. Returning that lower endpoint cannot enlarge the
failure event. At $k=0$, zero is valid directly. The certified event
implies actual gain at least $t$, so nonnegativity gives
$\mathbb E[G\mid\mathcal H_A]\ge tp$. A union bound over 232
events has cost $232\alpha_B=1/100$.

Adding the three stated family allocations gives $1/20$. Keeping
their actual Hoeffding expressions instead gives
$(464+696)\exp(-49/4)+1/100$, which is strictly below
$1160/200000+1/100=79/5000$. This proves both coverage statements.
The bound concerns all fixed cells, including zero-event, unchanged,
failed-enhancement and unfavorable cells; no selection by outcome is
permitted.
\end{proof}

\begin{corollary}[Simultaneous coverage of the unchanged first and new second families]
\label{cor:jointinference124ab}
Retain every first-experiment interval and threshold bound from
Proposition~\ref{prop:inference124}, without changing its radius,
threshold, family or observations. Combine them with the prospective
second-experiment family above. Under the two experiments' numerical
and sampling premises, all claims hold jointly with probability
strictly greater than $1-3631/125000>19/20$.
This is a joint repeated-sampling statement for the adaptive two-stage
design; it does not assert coverage conditional on a particular
realized first-experiment confidence event.
\end{corollary}

\begin{proof}
The first family has 116 nonnegative gains and 290 signed comparisons,
each with two-sided Hoeffding exponent $25/2$, plus its threshold-family
failure $1/100$. Thus its actual failure probability is at most
\[
 \varepsilon_A=(232+580)\exp(-25/2)+1/100.
\]
The same exact 31-term series used to justify its published, looser
allocation also proves
\[
 \sum_{j=0}^{30}\frac{(25/2)^j}{j!}>250000.
\]
This strengthens the bound on the unchanged confidence construction;
it does not narrow any already reported first-experiment interval.
Let $E_A,E_B$ denote failure of any claim in the respective fixed
families. The second theorem gives, uniformly in all admissible
pre-draw preparation,
$\Pr(E_B\mid\mathcal H_A)\le\varepsilon_B$.
The law of total probability and a union bound give
\begin{align*}
 \Pr(E_A\cup E_B)
 &\le\Pr(E_A)+\mathbb E[\Pr(E_B\mid\mathcal H_A)]\\
 &\le812\exp(-25/2)+1160\exp(-49/4)+2/100\\
 &<812/250000+1160/200000+2/100\\
 &=3631/125000=0.029048<1/20.
\end{align*}
The first inequality requires no independence between $E_A$ and $E_B$.
Only the new draws must be independent conditional on the preparation.
The new design may use the first experiment's diagnosed limitations,
provided the whole second family is fixed before its own draw.
No unpublished experiment or extra observation is supplied by this
argument.
\end{proof}

As before, the reported root-action RMS certificate follows from
the original $\mu$-strong convexity and the constrained minimizer's
variational inequality:
\[
 |a-a^*|^2\le\frac{2}{\mu}\{q(a)-q(a^*)\}
       \le\frac{2}{\mu}\{J(a)-q(a^*)\}.
\]
For $\mu=2$, average the sound per-state regret upper bounds first,
then take an upward-rounded square root. This is a finite-sample upper
bound on RMS distance to the optimal-future root action; it is not
the observed exact action error, nor a population RMS confidence
bound unless a separate such construction is specified. The new signed
radius $7\tau/8$ is smaller than the support half-width, so its sign
can in principle be informative. Neither the formula nor the new
sample size promises that an economically useful ordering or a positive
complete-cost adoption domain will be found.
